\begin{table*}[!t]
\centering\small\setlength\tabcolsep{2pt}\renewcommand{\arraystretch}{1.05}
\caption{What this corpus reports, and what it does not. Five conclusions, each with what the corpus reports,
the claim it stops short of, the property that limits it most, and the experiment that would
test it. Row labels refer to the statements of \S\ref{sec:findings}; counts are on all records of the corpus. This is the authors' synthesis of \S\ref{sec:methods}--\S\ref{sec:future}, not a coded field.}
\label{tab:synthesis}
\begin{tabular}{@{}L{2.2cm}L{3.9cm}L{3.4cm}L{3.1cm}L{3.2cm}@{}}
\toprule
Question & Reported in this corpus & Not established & Principal limitation & Experiment that would test it \\
\midrule
Availability against use & Clinical findings, anatomy, demographic and acquisition attributes are reported
recoverable from frozen encoders in at least one study in each of four to seven modalities (F1a--F1c, F2,
F2b) & That the output uses them:
\Vspca{}--\Vspc{} of the \Vspcn{} applicable intervention records report a matched control, none meets the
criterion of \S\ref{sec:methods} for a concept-specific causal contribution (F9a) & Breadth is breadth over records, and the safeguards pooled into
one control are not equivalent & Pair a probe with a matched intervention or erasure at the same locus and
behavioural end-point (\S\ref{sec:future}, item 2) \\
Acquisition and site structure & Centre, scanner and acquisition structure is recoverable from pathology and
brain-MRI embeddings (F2b, F3a, F3c) & That removing it restores generalisation: the downstream effect is
operation-specific, and post-hoc harmonisation raising invariance is an operation-specific observation, not
a general finding & F3b rests on five records from five studies, two of them flagged by the defect screen; pathology dominates & Report the downstream end-point, not only invariance, after
harmonisation on a held-out site \\
Meaning of discovered features & \Nsaeval{} of the \Nsae{} dictionary studies attempt semantic
confirmation of feature meaning against clinicians or structured labels, all encoder-side (F7) & Cross-seed identity of features or names, which no
study establishes; feature completeness or causal meaning, which semantic agreement does not show; and any
confirmation of a decoder-residual dictionary & The row selects the studies that
attempted validation, so it records how validation was done, not how often it succeeds & Train two
differently seeded dictionaries on one encoder; report matched-feature fraction and validate names against
region annotations (\S\ref{sec:future}, item 3) \\
Medical against general pre-training & In the same-data comparisons of this corpus, selected geometric and
shortcut properties are reported in both kinds of model (F10, F11) & An effect of medical pre-training:
better clinical decodability in the medical model appears only in uncontrolled comparisons & \Nmvgrec{} same-data comparisons
from \Nmvgstud{} studies vary architecture, objective, scale and data together; most strata hold one to
five records & Paired training over matched but non-identical corpora, with volume, modality mix, objective, architecture and compute held fixed; or a mixture design varying the medical fraction of a fixed-size corpus in pre-specified steps, read as a dose--response
\\
Where to measure & The connector output and text encoder of multimodal models are almost unmeasured as loci
of their own (\S\ref{sec:loci}, F8b) & That apparent decodability depends on the readout and the locus,
which three heterogeneous studies suggest but none tests directly; and any ordering of the contributions of
model, depth and pooling, which no design compares in one experiment & The readout evidence is three
studies varying different things, and the locus evidence one uncontrolled successive-locus probe of class
labels & Probe one finding across encoder, connector output and every language-model layer of one report
generator (\S\ref{sec:future}, item 1) \\
\bottomrule
\end{tabular}
\end{table*}
